\section{The original twisted ratio line on the same normalized model}

\begin{lemma}[Pulling back the same original line along an invariant projection]
\label{lem:pl-original-invariant-projection-line-comparison}
\lean{EnvelopeProperty.PLOriginalModelCompatibility.invariantProjectionLineComparison}
\leanok
Let $a,p:G\times N\to N$ be the actual action and projection, and let
$q:N\to X$ be the unchanged original projection. If $a q=p q$ as whole
scheme morphisms, every original $\mathcal O_X$-module $L$ has a canonical
whole comparison $a^*q^*L\simeq p^*q^*L$.
\end{lemma}
\begin{proof}
Use the canonical composition comparison to identify $a^*q^*L$ with
$(a q)^*L$, the pullback congruence for the actual equality $a q=p q$,
and the inverse composition comparison to reach $p^*q^*L$.
\end{proof}

\begin{theorem}[Whole group comparison of the original normalized polarization]
\label{thm:pl-original-normalized-twisted-ratio-group-comparison}
\lean{EnvelopeProperty.ProjSubvariety.normalizedSimultaneousOriginalTwistNativeGmComparison,EnvelopeProperty.ProjSubvariety.normalizedSimultaneousTwistedRatioNativeGmComparison}
\leanok
\uses{thm:pl-original-normalized-projection-action-invariant,lem:pl-original-invariant-projection-line-comparison,thm:pl-original-normalized-cartier-factor-linearizations}
The same actual original normalized PL model, with the same actual group
action, carries an explicit whole group comparison on its original line
$q^*\mathcal O_X(n)\otimes(I_N\otimes J_N^\vee)$. Its input consists only
of the original homogeneous ideals, their actual coefficient-unit
invariance and finite Gauss logs; no action on this line or comparison of
the resulting modules is assumed.
\end{theorem}
\begin{proof}
Apply the preceding lemma to the original twist line on the original
variety. The original ratio factors are actual invertible ideals on the
same normalized model, invariant under the actual group action. Their
constructed canonical ratio comparison applies. Compose the
canonical tensor pullback comparison, the tensor product of these two
whole comparisons, and the inverse tensor pullback comparison for the
projection. This is a comparison on the unchanged original tensor line
over the entire original group product.
\end{proof}
